\begin{lemma}[One round of the Schauder bootstrap]
Suppose interior divergence-form Schauder estimates hold, in the sense of
\noderef{DivFormSchauderInput}.  Let $j\in\mathbb{N}$, let $\kappa,\lambda,\Lambda,M,b,r,r'$ be
real numbers with
\[
  0<\kappa<1,\qquad 0<\lambda\le\Lambda,\qquad M>0,\qquad b>0,\qquad 0<r'<r .
\]
Then there is $N>0$, depending only on $j,\kappa,\lambda,\Lambda,M,b,r,r'$, with the following
property.

Let $v\colon\mathrm{Plane}\to\mathbb{R}$ and
$A\colon\mathrm{Plane}\to\mathcal{L}(\mathrm{Plane},\mathrm{Plane})$, write
$Dv\colon\mathrm{Plane}\to\mathrm{Plane}$ for the classical gradient field of $v$ and
$\partial_ev(x):=\langle Dv(x),e\rangle$ for $e\in\mathrm{Plane}$, and assume
\begin{enumerate}
\item[(H1)] $v$ is of class $C^{\infty}$ on $B_r$;
\item[(H2)] $\lvert Dv(x)\rvert\le b$ for every $x\in B_r$;
\item[(H3)] for every $x\in B_r$ and every $\xi\in\mathrm{Plane}$,
  $\lambda\lvert\xi\rvert^{2}\le\langle A(x)\xi,\xi\rangle$ and
  $\lvert A(x)\xi\rvert\le\Lambda\lvert\xi\rvert$;
\item[(H4)] $A$ is of class $C^{j}$ on $B_r$ and
  $\lVert D^{i}A(x)\rVert\le M$ for every $i\le j$ and every $x\in B_r$;
\item[(H5)] $\lVert D^{j}A(x)-D^{j}A(y)\rVert\le M\lvert x-y\rvert^{\kappa}$ for all
  $x,y\in B_r$;
\item[(H6)] for every $e\in\mathrm{Plane}$ the quadruple
  $\bigl(B_r,\;A,\;\partial_ev,\;D(\partial_ev)\bigr)$ satisfies the five conditions of
  \noderef{IsLinearDivSol}.
\end{enumerate}
Then
\[
  \bigl\lVert D^{m}v(x)\bigr\rVert\le N
  \quad(1\le m\le j+2,\ x\in B_{r'}),
  \qquad
  \bigl\lVert D^{\,j+2}v(x)-D^{\,j+2}v(y)\bigr\rVert\le N\lvert x-y\rvert^{\kappa}
  \quad(x,y\in B_{r'}).
\]

In words: one round of the bootstrap converts $C^{j,\kappa}$ bounds on the coefficient field
$A$, together with the differentiated equations (H6), into $C^{\,j+2,\kappa}$ bounds on the
scalar function $v$ on a smaller ball.  The gain of two orders over (H4) is the Schauder gain of
one order applied to the first derivatives of $v$.  The conclusion begins at order $1$: a bound
at order $0$ is a bound on $v$ itself, for which no Schauder estimate is needed, and none is
asserted here.  The constant $N$ is produced before $v$ and $A$, so the round is uniform over
all data satisfying the stated bounds.
\end{lemma}
\begin{proof}
Let $(e_1,e_2)$ be the standard orthonormal basis of $\mathrm{Plane}$, so $\lvert e_l\rvert=1$
and every $z\in\mathrm{Plane}$ satisfies
\begin{equation}\label{eq:br-split}
  z=\sum_{l=1}^{2}\langle z,e_l\rangle\,e_l .
\end{equation}

\medskip\noindent\emph{A remark on post-composition with a continuous linear map, used twice
below.}
\begin{itemize}
\item[(L)] Let $Z,Z'$ be normed real vector spaces and let $L\colon Z\to Z'$ be linear and
continuous.  Then $L$ is Fr\'echet differentiable at every $\zeta\in Z$ with $DL(\zeta)=L$:
linearity gives $L(\zeta+\eta)-L(\zeta)-L(\eta)=0$ for every $\eta$, so the remainder in the
definition of the Fr\'echet derivative vanishes identically.  Consequently, by clause (1) of the
background proposition on the chain rule for compositions, if $G\colon\mathrm{Plane}\to Z$ is
differentiable at $x$ then so is $L\circ G$, with $D(L\circ G)(x)=L\circ DG(x)$.  Moreover, if
$U_0\subseteq\mathrm{Plane}$ is open and $G$ is of class $C^{\infty}$ on $U_0$, then for every
$n\in\mathbb{N}$, every $x\in U_0$ and all $u_1,\ldots,u_n\in\mathrm{Plane}$,
\begin{equation}\label{eq:br-L}
  D^{n}(L\circ G)(x)[u_1,\ldots,u_n]=L\bigl(D^{n}G(x)[u_1,\ldots,u_n]\bigr),
\end{equation}
equivalently $D^{n}(L\circ G)(x)=\Lambda_n\bigl(D^{n}G(x)\bigr)$, where
$\Lambda_n\colon Q\mapsto\bigl(u_1,\ldots,u_n\mapsto L(Q[u_1,\ldots,u_n])\bigr)$ carries
continuous $n$-linear maps $\mathrm{Plane}^{n}\to Z$ to continuous $n$-linear maps
$\mathrm{Plane}^{n}\to Z'$.  The map $\Lambda_n$ is linear, and continuous because
$\lVert\Lambda_n(Q)\rVert\le\lVert L\rVert\lVert Q\rVert$.  The identity
\eqref{eq:br-L} is proved by induction on $n$: for $n=0$ both sides are $L(G(x))$ by clause (1)
of the background proposition on the recursion for the iterated Fr\'echet derivative; and for the
step, $z\mapsto D^{n}G(z)$ is differentiable on $U_0$ by clause (3) of the recursion proposition,
so the first part of (L) applied to the continuous linear map $\Lambda_n$ gives
$D\bigl(D^{n}(L\circ G)\bigr)(x)[w]=\Lambda_n\bigl(D(D^{n}G)(x)[w]\bigr)$, and evaluating at
$(u_1,\ldots,u_n)$ and using clause (2) of the recursion proposition on each side yields the
identity at $n+1$.

Finally, if $L$ is a linear \emph{isometry} then so is each $\Lambda_n$: it is linear, and
\[
  \lVert\Lambda_n(Q)\rVert
  =\sup\{\lVert L(Q[u_1,\ldots,u_n])\rVert:\lvert u_1\rvert\le1,\ldots,\lvert u_n\rvert\le1\}
  =\sup\{\lVert Q[u_1,\ldots,u_n]\rVert:\lvert u_1\rvert\le1,\ldots,\lvert u_n\rvert\le1\}
  =\lVert Q\rVert ,
\]
the outer equalities by the background proposition on the operator norm of a continuous
multilinear map.
\end{itemize}

\medskip\noindent\emph{Step 1: the constants.}
Put $c:=b+1>0$, so $b\le c$.  Apply \noderef{DivFormSchauderInput} with the parameters
$j,\kappa,\lambda,\Lambda,M,r,r'$ --- admissible since $0<\kappa<1$, $0<\lambda\le\Lambda$,
$M>0$ and $0<r'<r$ --- and let $S>0$ be the constant it furnishes; it depends only on
$j,\kappa,\lambda,\Lambda,M,r,r'$.  Put
\[
  N:=2cS>0 ,
\]
which depends only on $j,\kappa,\lambda,\Lambda,M,b,r,r'$, hence was fixed before $v$ and $A$.

\medskip\noindent\emph{Step 2: the gradient field, and the bridge to $v$.}
By (H1) and clause (1) of the background proposition on regularity of the gradient field and of
directional derivatives, applied with $n=\infty$, the field $Dv$ is of class $C^{\infty}$ on
$B_r$; hence so is each $\partial_{e_l}v=\langle Dv(\cdot),e_l\rangle$, by clause (2) of that
proposition.  Write $D^{1}v\colon\mathrm{Plane}\to\mathrm{Plane}^{*}$ for the first Fr\'echet
derivative of $v$.  Apply \noderef{GradientFieldOrderShift} with the open set $U:=B_r$ --- open
as a metric ball --- and the function $v$, which is of class $C^{\infty}$ there by (H1).  Its
clauses (3) and (4) are the \emph{bridge}
\begin{equation}\label{eq:br-bridge}
  \bigl\lVert D^{n}(Dv)(x)\bigr\rVert=\bigl\lVert D^{\,n+1}v(x)\bigr\rVert,
  \qquad
  \bigl\lVert D^{n}(Dv)(x)-D^{n}(Dv)(y)\bigr\rVert
    =\bigl\lVert D^{\,n+1}v(x)-D^{\,n+1}v(y)\bigr\rVert
\end{equation}
for every $n\in\mathbb{N}$ and all $x,y\in B_r$.  Both are equalities: nothing is lost in passing
between the scalar function and its gradient field.

\medskip\noindent\emph{Step 3: the reassembly.}
For $l\in\{1,2\}$ let $\iota_l\colon\mathbb{R}\to\mathrm{Plane}$, $\iota_l(t):=t\,e_l$.  It is
linear, and isometric because $\lvert t\,e_l\rvert=\lvert t\rvert\lvert e_l\rvert=\lvert t\rvert$.
Put $h_l:=\iota_l\circ\bigl(\partial_{e_l}v\bigr)$, so that $h_l(z)=\bigl(\partial_{e_l}v(z)\bigr)e_l$,
and by \eqref{eq:br-split} applied to $z:=Dv(z)$,
\begin{equation}\label{eq:br-comp}
  Dv=h_1+h_2\qquad\text{on }\mathrm{Plane}.
\end{equation}
Each $\partial_{e_l}v$ is of class $C^{\infty}$ on $B_r$ by Step 2, hence so is $h_l$, by (L)
and the continuity and linearity of $\iota_l$.  Since $D^{n}$ is additive in the map --- clause
(2) of the recursion proposition together with the linearity of the Fr\'echet derivative give
$D^{n}(f+g)=D^{n}f+D^{n}g$ for maps of class $C^{n}$ --- \eqref{eq:br-comp} gives
$D^{n}(Dv)(x)=D^{n}h_1(x)+D^{n}h_2(x)$ for $x\in B_r$.  By (L) with $L:=\iota_l$, which is an
isometry, $\lVert D^{n}h_l(x)\rVert=\lVert D^{n}(\partial_{e_l}v)(x)\rVert$, and likewise, the
corresponding $\Lambda_n$ being linear,
$\lVert D^{n}h_l(x)-D^{n}h_l(y)\rVert=\lVert D^{n}(\partial_{e_l}v)(x)-D^{n}(\partial_{e_l}v)(y)\rVert$.
With the triangle inequality and \eqref{eq:br-bridge} we conclude, for every $n\in\mathbb{N}$ and
all $x,y\in B_r$,
\begin{equation}\label{eq:br-reassemble}
  \bigl\lVert D^{\,n+1}v(x)\bigr\rVert\le\sum_{l=1}^{2}
    \bigl\lVert D^{n}(\partial_{e_l}v)(x)\bigr\rVert,
  \qquad
  \bigl\lVert D^{\,n+1}v(x)-D^{\,n+1}v(y)\bigr\rVert\le\sum_{l=1}^{2}
    \bigl\lVert D^{n}(\partial_{e_l}v)(x)-D^{n}(\partial_{e_l}v)(y)\bigr\rVert .
\end{equation}

\medskip\noindent\emph{Step 4: normalizing each directional derivative.}
Fix $l\in\{1,2\}$ and put $w:=\partial_{e_l}v$, $G:=Dw$.  By (H6) with $e:=e_l$, the quadruple
$(B_r,A,w,G)$ satisfies the five conditions of \noderef{IsLinearDivSol}.  So does
$(B_r,A,c^{-1}w,c^{-1}G)$.  Indeed condition (1) holds because a scalar multiple of a function
continuous on $B_r$ is continuous on $B_r$; condition (2) because a scalar multiple of an almost
everywhere strongly measurable map is almost everywhere strongly measurable; condition (3)
because $\lvert c^{-1}G\rvert^{2}=c^{-2}\lvert G\rvert^{2}$ is integrable on a compact subset of
$B_r$ whenever $\lvert G\rvert^{2}$ is, a constant multiple of an integrable function being
integrable; condition (4) because both of its integrals are homogeneous of degree one in the
pair $(w,G)$, so multiplying the identity for $(w,G)$ by $c^{-1}$ gives the identity for
$(c^{-1}w,c^{-1}G)$; and condition (5) because $A(x)$ is linear, so
$A(x)\bigl(c^{-1}G(x)\bigr)=c^{-1}A(x)\bigl(G(x)\bigr)$, the inner product is linear in its first
argument, and the integral is therefore again homogeneous of degree one, so that it vanishes for
$c^{-1}G$ as it does for $G$.

Moreover $c^{-1}w$ is normalized: for $x\in B_r$, by the Cauchy--Schwarz inequality for the inner
product of \noderef{Plane}, $\lvert e_l\rvert=1$ and (H2),
\[
  \bigl\lvert c^{-1}w(x)\bigr\rvert
  =c^{-1}\bigl\lvert\langle Dv(x),e_l\rangle\bigr\rvert
  \le c^{-1}\lvert Dv(x)\rvert
  \le c^{-1}b\le1 ,
\]
the last step because $b\le c$ and $c>0$.

\medskip\noindent\emph{Step 5: one Schauder round, and undoing the scaling.}
Apply \noderef{DivFormSchauderInput} with the parameters of Step 1 and with the coefficient field
$A$, the solution $c^{-1}w$ and its weak gradient $c^{-1}G$.  Its three hypotheses hold: uniform
ellipticity on $B_r$ is (H3); the $C^{j,\kappa}$ bounds on $A$ are (H4) and (H5); and the
normalized-weak-solution hypothesis is Step 4.  Its conclusion gives
\[
  \bigl\lVert D^{m}\bigl(c^{-1}w\bigr)(x)\bigr\rVert\le S\ \ (m\le j+1,\ x\in B_{r'}),
  \quad
  \bigl\lVert D^{\,j+1}\bigl(c^{-1}w\bigr)(x)-D^{\,j+1}\bigl(c^{-1}w\bigr)(y)\bigr\rVert
    \le S\lvert x-y\rvert^{\kappa}
\]
for $x,y\in B_{r'}$.  Now $D^{m}$ is linear in the map, by clause (2) of the recursion
proposition together with the linearity of the Fr\'echet derivative, so
$D^{m}\bigl(c^{-1}w\bigr)=c^{-1}D^{m}w$ for every $m$; multiplying through by $c>0$ gives, for
each $l\in\{1,2\}$,
\begin{equation}\label{eq:br-w}
  \bigl\lVert D^{m}(\partial_{e_l}v)(x)\bigr\rVert\le cS\ \ (m\le j+1,\ x\in B_{r'}),
  \quad
  \bigl\lVert D^{\,j+1}(\partial_{e_l}v)(x)-D^{\,j+1}(\partial_{e_l}v)(y)\bigr\rVert
    \le cS\lvert x-y\rvert^{\kappa}
\end{equation}
for $x,y\in B_{r'}$.

\medskip\noindent\emph{Conclusion.}
Note $B_{r'}\subseteq B_r$, since $r'<r$.  Let $1\le m\le j+2$ and $x\in B_{r'}$, and write
$m=n+1$ with $n:=m-1\le j+1$.  The first half of \eqref{eq:br-reassemble}, applied at $x\in B_r$,
together with the first half of \eqref{eq:br-w} at the order $n\le j+1$, gives
\[
  \bigl\lVert D^{m}v(x)\bigr\rVert\le\sum_{l=1}^{2}
    \bigl\lVert D^{n}(\partial_{e_l}v)(x)\bigr\rVert\le2cS=N .
\]
For the H\"older bound, let $x,y\in B_{r'}$ and take $n:=j+1$ in the second half of
\eqref{eq:br-reassemble}; with the second half of \eqref{eq:br-w},
\[
  \bigl\lVert D^{\,j+2}v(x)-D^{\,j+2}v(y)\bigr\rVert
  \le\sum_{l=1}^{2}
    \bigl\lVert D^{\,j+1}(\partial_{e_l}v)(x)-D^{\,j+1}(\partial_{e_l}v)(y)\bigr\rVert
  \le2cS\lvert x-y\rvert^{\kappa}=N\lvert x-y\rvert^{\kappa} .
\]
These are the two assertions.
\end{proof}
